\fancyhead[L]{ARSE41}
\renewcommand{\arraystretch}{1}

\titre{}{ARSE41 - Interface Homme Machine}


\cadreR{\defi  Nous allons nous intéresser aux \textbf{interfaces Homme-Machine (IHM)}. Une IHM permet à un utilisateur d'interagir avec une machine, la souris et le clavier sont des exemples d'IHM, mais il en existe bien d'autres (par exemple les écrans tactiles des smartphones). Par "machine", on entend bien évidemment "ordinateur", comme dans le cas du duo clavier-souris, mais pas seulement. En effet, les "machines" peuvent aussi être des objets connectés, des systèmes embarqués ou encore des robots.  }


\section{La carte Arduino}


\cadreR{\defi \textbf{Arduino} est un \textbf{circuit imprimé} en matériel libre sur lequel se trouve un \textbf{microcontrôleur} qui peut être programmé pour analyser et produire des \textbf{signaux électriques}.}

\medskip

Cette carte possède:

\begin{liste}
	\item un \textbf{microcontroleur}: C'est le cœur de la carte Arduino. C'est une sorte d'ordinateur minuscule (mémoire morte, mémoire vive, processeur et entrées/sorties) et c'est lui que nous allons programmer.
	\item des \textbf{périphériques d'entrée}
	\item des \textbf{périphériques de sortie}
\end{liste}


\begin{center}
\includegraphics[scale=0.4]{ScreenShot_20240705085231.ps}
\end{center}


\begin{center}
\includegraphics[scale=0.35]{ScreenShot_20240705085524.ps}
\end{center}

\newpage
\section{Programmation de la carte Arduino}

La programmation se fait dans un langage proche du langage \verb+C+. Son fonctionnement est le suivant:

\begin{center}
\includegraphics[scale=0.4]{ScreenShot_20240705090052.ps}
\end{center}


\medskip

La compréhension du code est plutôt simple. en voici un exemple.

\begin{center}
\includegraphics[scale=0.35]{ScreenShot_20240705090207.ps}
\end{center}

\medskip

Une fois le programme tapé dans ce langage, il faut ensuite le \textbf{compiler} de manière à ce qu'il soit compréhensible par la machine

\begin{center}
\includegraphics[scale=0.6]{2575329f-0088-4c95-8a8e-da4bfc1d8c06.png.960x960_q85.ps}
\end{center}
\import{../../../../preambule}{raz_compteurs.tex}
\newpage
